%% kapa: exam=1 round=22 year=2011 period=2 form=A subject=accounting subject_ko=회계학 seq=23 number=23
%% source: Q-Net 감정평가사 기출문제 (cst00309) artlSeq=5036790 file=extracted/1008/5036790/2011년 제22회 감정평가사 1차시험 기출문제/2교시.hwp
%% conversion: kordoc 4.24.51 -> tools/gen1.py (auto); review=proofread-gpt-astra; flags=none
\begin{question}{23}
(주)서울은 20×1년 초에 기계장치(취득원가 ￦100,000, 감가상각누계액 ￦20,000)를 다음과 같은 조건(ㄱ, ㄴ, ㄷ) 가운데 하나로 (주)한국의 유형자산과 교환하였다. (주)서울의 입장에서 유형자산처분이익이 높은 순서로 배열된 것은? (단, 각 거래는 독립적인 상황으로 가정함)
\begin{framedtext}
ㄱ. (주)서울의 기계장치 공정가치는 ￦85,000이며, (주)한국의 건물과 교환하였다. (주)서울은 교환시 현금 ￦15,000을 지급하였다. 단, 이 거래는 상업적실질이 존재하는 거래이다. \par
ㄴ. (주)서울의 기계장치 공정가치는 ￦90,000이며, (주)한국의 토지와 교환하였다. (주)서울은 교환시 현금 ￦20,000을 수령하였다. 단, 이 거래는 상업적실질이 존재하는 거래이다. \par
ㄷ. (주)서울의 기계장치 공정가치는 ￦90,000이며, (주)한국의 동종 기계장치와 교환하였다. (주)서울은 교환시 현금 ￦25,000을 수령하였다. 단, 이 거래는 상업적실질이 존재하지 않는 거래이다. \par
\end{framedtext}
\begin{choices}
  \item ㄱ > ㄴ > ㄷ
  \item ㄱ > ㄷ > ㄴ
  \item ㄴ > ㄱ > ㄷ
  \item ㄴ > ㄷ > ㄱ
  \item ㄷ > ㄴ > ㄱ
\end{choices}
\end{question}
\answer{3}
